% ============================================================
%  H. Høffding, "Om Realisme i Videnskab og Tro"
%  1884
%  TRANSLATION (English)
% ============================================================
\documentclass[12pt,a4paper]{article}
\usepackage[utf8]{inputenc}
\usepackage[T1]{fontenc}
\usepackage[english]{babel}
\usepackage{microtype}
\usepackage{setspace}
\usepackage[margin=1.2in]{geometry}
\usepackage{fancyhdr}
\usepackage{hyperref}

\hypersetup{
  pdftitle={On Realism in Science and Faith},
  pdfauthor={H. Høffding},
  hidelinks
}

\pagestyle{fancy}
\fancyhf{}
\rhead{Høffding, \emph{On Realism in Science and Faith} (1884)}
\cfoot{\thepage}

\setstretch{1.3}

\title{\textbf{On Realism in Science and Faith}}
\author{H.\ Høffding\\[4pt]
  \small Translated from the Danish}
\date{1884}

\begin{document}
\maketitle
\thispagestyle{fancy}

\bigskip

\noindent The word \emph{realism} is nowadays most frequently used to
designate a tendency in poetic literature. It is sometimes employed
in a reproachful sense, sometimes in a laudatory one, but rarely with
a clear awareness of what it actually expresses, or of the limits of
the validity of the concept it names. A clear account of this may
perhaps only emerge once the modern aesthetic movement has more fully
and distinctly unfolded its character and its consequences. It is not
aesthetic realism that I intend to discuss here. Rather, I shall
attempt to bring out what is distinctive about the realist tendency
that, in our time, manifests itself no less clearly in the domains of
science and of general outlook on life. It follows of itself that
there must be an inner connection among the tendencies at work in
these different domains, for just as little in the spiritual world as
in the physical is there anything accidental or isolated. The human
mind cannot be partitioned in such a way that, in its imagination, it
follows paths that have absolutely nothing analogous or parallel to
those it takes in its thinking and in its faith. Only the most
doctrinaire and rigid among the opponents of modern aesthetic realism
assign art and poetry such a sharply segregated place in the life of
the mind that currents from other domains cannot---or ought not---reach
them. By examining the nature and justification of realism in one
domain, we indirectly contribute to the understanding of its
appearance in others.

\section*{I.}

In the most comprehensive sense, we may define realism as the principle
of natural causes. If realism has opponents, and if it has only slowly
worked its way forward to clear self-consciousness, this must be
because the explanation of phenomena and events has been sought---and
is still being sought---by routes other than those of natural causes.

\medskip

\noindent The human drive to find the causes of things has expressed
itself long before the founding of a realistic science. No period can
be pointed to in the developmental history of the human mind in which
it stood entirely passive before sense-impressions, without the
slightest attempt to explain them in some way or other.

\medskip

\noindent The contrast between different periods in the history of
cognition is instead characterized by the different explanatory or
interpretive principles applied to experience. And here, above all,
two great principles stand in sharp opposition to one another.

The one interprets natural phenomena according to rules drawn from
outside nature; the other interprets nature through nature itself.

\medskip

\noindent Under the first kind of natural explanation belongs all
mythology. The mythological conception of nature explains phenomena
as expressions of the will of personal beings. For it, nature is
populated by more or less human-like spirits, whose intervention is
traced above all in remarkable and significant events.

\medskip

\noindent Something similar holds as well of many forms of philosophical
idealism. If one finds the true explanation of a natural phenomenon
in the purpose it is assumed to serve, or in the idea it is supposed
to express, one is not explaining nature through nature itself.

\medskip

\noindent Realism, with the principle of natural causes, always derives
natural phenomena from other natural phenomena, the more complex
phenomena from the simpler. When it is said that nature must be
explained through itself, it must be remembered that nature is an
immeasurable world of phenomena unfolding for us in time and in
space. The principle that nature is to be explained through itself
therefore sets an infinite task.

\section*{II.}

When we think through the concept of realism in this way, the
absolutely hostile opposition to idealism falls away. It is one of
idealism's fundamental thoughts that we cannot rest at the level of
individual, isolated facts, but must seek a bond that can bring
coherence and unity to bear. Idealism went wrong in assuming that
we could reach this unity and coherence by any route other than the
continued working-through of the facts of experience.

\medskip

\noindent Is realism, then, proven? Not at all. And it can be proven
in the strict sense just as little as any principle can be. We arrive
at a principle when, by following a developmental course, we discover
a dominant rule or tendency that becomes more clearly visible the
further the development proceeds. Every principle is, to that extent,
a hypothesis that we try to apply to reality.

\section*{III.}

The expression ``natural causes'' requires closer explanation. If by
the natural we understand that which can become an object of secure
human experience, then spiritual phenomena belong to nature just as
much as material ones. Nature accordingly encompasses two groups or
series of phenomena, which we can bring into relation with one another
only by way of hypothesis or speculation: the material phenomena
unfolding in space, and the phenomena of conscious life, which become
immediately accessible to us only through self-observation, and which
can only be represented in spatial form in a symbolic way.

\medskip

\noindent The difficulty lies, as already noted, in finding the
connection between the two domains of experience. Materialism cuts
through the knot by overlooking their difference and simply making
the spiritual a form or product of the material.

\medskip

\noindent There is probably no other way out than to conceive of
consciousness-phenomena as inner, psychical expressions of the same
thing that, for the outer senses, appears as a material phenomenon.
The real existence is thus only one, but we apprehend it (at certain
points at least) in a double form.

\medskip

\noindent Realism therefore need not contest what is the genuine motive
and deepest interest of the idealistic worldview. It is quite
compatible with an effort to uphold the value and significance of
spiritual life in existence.

\section*{IV.}

Already in the last reflections we have been led from the domain of
science over into that of faith. For the most common distinguishing
mark between science and faith is that the former investigates the
laws and inner coherence of existence, while the latter asks about
the ethical significance and value of existence.

\medskip

\noindent When we define the concept of faith in this way, no necessary
conflict exists between science and faith. Science grounds and explains;
faith evaluates. The common foundation of both is the reality accessible
to our experience. Its laws are sought by science, supported by the
principle of natural causes, and faith expresses the value that this
law-governed reality must have for us in virtue of our ideals.

\medskip

\noindent By realism in faith I mean a view of life that holds fast to
the belief in the ideal value of world-development, while at the same
time being equally deeply convinced that this ideal value is realized
through the operation of natural causes.

\medskip

\noindent The customary meaning of the word faith, by contrast,
excludes such a harmony. It demands a more or less supernatural
apparatus for the realization of its ends, and sets up fixed forms
and inviolable symbols. On this rests the difference between humane
or philosophical faith and theological faith.

\medskip

\noindent Only \emph{blas\'etude} throws away the kernel along with
the husks. Such blas\'etude is surely the most dangerous enemy of
spiritual development in our age, more dangerous than all intolerance
and fanaticism, for in the latter there is at least life, while
blas\'etude is itself petrified and petrifies everything around it.

\end{document}
